\documentclass{article}
\usepackage[UTF8]{ctex}
\usepackage{amsmath,amssymb,amsfonts}
\usepackage{enumitem}
\usepackage{geometry}
\geometry{a4paper, margin=1in}
% === 错题本排版优化（不参与题目标识） ===
\setlength{\parindent}{0pt}
\setlength{\parskip}{0.35em plus 0.10em minus 0.08em}
\setlength{\emergencystretch}{3em}
\setlist{itemsep=0.25em, topsep=0.40em, parsep=0pt}
\allowdisplaybreaks[2]
% === 错题本排版优化结束 ===
\begin{document}
\section*{安徽大学2025年数学分析真题}
\section*{111．安徽大学2025年数学分析试卷}

\subsection{题目1}
一．（16 分）计算题．

\begin{itemize}
  \item[（1）] （8 分）设 \(a_{n}=\sum_{k=1}^{n}\left(\sqrt{1+\frac{k}{n^{2}}}-1\right)\) ，求 \(\lim _{n \rightarrow \infty} a_{n}\) 。
  \item[（2）] （8 分）求 \(\lim _{x \rightarrow 0^{+}} \frac{\int_{0}^{x} e^{-t^{2}} \mathrm{~d} t-x}{\sin x-x}\) ．
\end{itemize}

\subsection{题目2}
二、（14 分）证明 \(f(x, y)=\left\{\begin{array}{ll}\frac{x y}{\sqrt{x^{2}+4 y^{2}}}, & x^{2}+y^{2} \neq 0 ; \\ 0, & x^{2}+y^{2}=0 .\end{array}\right.\) 在原点 \((0,0)\) 处连续且可偏导，但不可微．

\subsection{题目3}
三、（15 分）计算

\[
\iint_{\Sigma}(x-y) \mathrm{d} y \mathrm{~d} z+(y-z) \mathrm{d} z \mathrm{~d} x+(z-x) \mathrm{d} x \mathrm{~d} y
\]

其中 \(\sum\) 为曲面 \(z=x^{2}+y^{2}(z \leq 1)\) ，方向向下。

\subsection{题目4}
四．（15 分）讨论反常积分 \(\int_{0}^{+\infty} \frac{e^{\sin x} \sin 2 x}{x^{p}} \mathrm{~d} x\) 的敛散性，其中 \(p>0\) ．

\subsection{题目5}
五．（15 分）设 \(f(x)=\sum_{n=0}^{\infty} \frac{1}{2^{n}+x}\) ．\\
（1）（10 分）证明函数 \(f(x)\) 在 \([0,+\infty)\) 上可导且一致连续．\\
（2）（5 分）证明反常积分 \(\int_{0}^{+\infty} f(x) \mathrm{d} x\) 发散．

\subsection{题目6}
六．（15 分）若函数 \(f\) 一阶导数或偏导数存在且连续，则记 \(f \in C^{1}\) ．设函数 \(z=f(x, y)\) 由方程

\[
F(x z-y, x-y z)=0
\]

唯一确定，其中 \(f \in C^{1}, F \in C^{1}\) ，试求 \(\oint_{L}\left(x z^{2}+2 y z\right) \mathrm{d} y-\left(2 x z+y z^{2}\right) \mathrm{d} x\) ，其中 \(L\) 为正向单位圆周．

\subsection{题目7}
七．（15 分）若将 \(a_{1}, a_{2}, \cdots, a_{n}, \cdots\) 重新排列后得到 \(a_{1}^{\prime}, a_{2}^{\prime}, \cdots, a_{n}^{\prime}, \cdots\) ，则级数 \(\sum_{n=1}^{\infty} a_{n}^{\prime}\) 称为 \(\sum_{n=1}^{\infty} a_{n}\) 的一个更序级数。设级数 \(\sum_{n=1}^{\infty} x_{n}\) 条件收敛，证明：对任意给定的 \(a,-\infty<a<+\infty\) ，必定存在 \(\sum_{n=1}^{\infty} x_{n}\) 的更序级数 \(\sum_{n=1}^{\infty} x_{n}^{\prime}\) ，满足 \(\sum_{n=1}^{\infty} x_{n}^{\prime}=a\) 。

\subsection{题目8}
八．（15 分）设 \(f(x)\) 在 \([0,+\infty)\) 上连续且有界，若 \(\forall r \in(-\infty,+\infty), f(x)=r\) 在 \([0,+\infty)\) 上只有有限个根或无根，证明： \(\lim _{x \rightarrow+\infty} f(x)\) 存在．

\subsection{题目9}
九．（15 分）设有界函数 \(f(x)\) 在闭区间 \([a, b]\) 上的不连续点为 \(\left\{x_{n}\right\}_{n=1}^{\infty}\) ，且 \(\lim _{n \rightarrow \infty} x_{n}\) 存在，证明：\(f(x)\)在闭区间 \([a, b]\) 上可积．

附加题．（15 分）设 \(f(x)\) 在闭区间 \([a, b]\) 上二阶可导，且 \(f(a)=f(b)=0\) ．证明：对任意的 \(x \in(a, b)\) ，存在 \(\xi \in(a, b)\) ，使得

\[
f(x)=\frac{f^{\prime \prime}(\xi)(x-a)(x-b)}{2} .
\]
\end{document}
